%%% Verze pro jednostranný tisk:
\documentclass[11pt,a4paper]{report}
\usepackage[top=25mm,bottom=25mm,right=25mm,left=30mm,head=12.5mm,foot=12.5mm]{geometry}
\let\openright=\clearpage

%%% Pokud tiskneme oboustranně:
%\documentclass[11pt,a4paper,twoside,openright]{report}
%\usepackage[top=25mm,bottom=25mm,right=25mm,left=30mm,head=12.5mm,foot=12.5mm]{geometry}
%\let\openright=\cleardoublepage

%%% DEFINICE ZÁKLADNÍCH PROMĚNNÝCH
\def\TypPrace{BP}                % bakalářská práce/bachelor thesis
%\def\TypPrace{DP}               % diplomová práce/master thesis
\def\Jazyk{cze}                  % čeština/czech
%\def\Jazyk{slo}                 % slovenština/slovak
%\def\Jazyk{eng}                 % angličtina/english

%%% Definice různých užitečných maker (viz popis uvnitř souboru)
\input{macros}

%%% Název práce v jazyce práce (přesně podle zadání)
%%% Title of the thesis in the language used in the text (exact according to assignment)
\def\NazevPrace{Návrh, implementace a evaluace MCP serveru pro zpřístupnění studijních dat VŠE AI agentům}

%%% Jméno autora
%%% Author's name - First name Surname
\def\AutorPrace{Jan Komínek}

%%% Rok odevzdání a měsíc (slovně)
%%% Year of submission and month (verbally) - month YYYY
\def\DatumOdevzdani{květen 2027}

%%% Vedoucí práce: Jméno a příjmení s~tituly
%%% Supervisor: First name and surname with titles
\def\Vedouci{Ing. Luboš Pavlíček}   % TODO ověřit přesné tituly vedoucího dle InSIS

%%% Konzultant práce: Jméno a příjmení s~tituly
%%% Consultant: First name and surname with titles
\def\Konzultant{}

%%% Studijní program
%%% Study program
\def\StudijniProgram{Aplikovaná informatika}   % TODO ověřit přesný název programu dle InSIS

%%% Studijní program - specializace
%%% Study program - specialization
\def\Specializace{}   % TODO doplnit specializaci, pokud ji program má

%%% Nepovinné poděkování (vedoucímu práce, konzultantovi, tomu, kdo zapůjčil software, literaturu apod.)
%%% Optional thanks (the supervisor, the consultant, the borrower of software, literature, etc.)
\def\Podekovani{%
Poděkování. % TODO
}

%%% Abstrakt (doporučený rozsah cca 150-250 slov; nejedná se o zadání práce)
\def\Abstrakt{%
Abstrakt. % TODO napsat na konci (cca 150–250 slov)
}
\def\AbstraktEN{%
Abstract. % TODO English abstract (~150–250 words)
}

%%% 3 až 5 klíčových slov (doporučeno)
\def\KlicovaSlova{Model Context Protocol, MCP server, velké jazykové modely, integrace dat, InSIS, studijní data, AI agenti}
\def\KlicovaSlovaEN{Model Context Protocol, MCP server, large language models, data integration, InSIS, study data, AI agents}

%%% Titulní strana a různé povinné informativní strany
%%% Title page and various mandatory information pages
\begin{document}
\include{beginning}

%%% Strana s automaticky generovaným obsahem práce
%%% A page with automatically generated content of the thesis
\setcounter{tocdepth}{2}
\tableofcontents

%%% Seznamy obrázků / tabulek / výpisů kódu zapneme, až jich bude v práci >20.
%\openright
%\listoffigures
%\clearpage
%\listoftables
%\clearpage
%\lstlistoflistings

%%% Seznam zkratek v práci (volitelně)
%%% List of abbreviations in the thesis (optionally)
\include{abbreviations}

\pagestyle{fancyx}

%%% ÚVOD (nečíslovaná kapitola)
{%
\pagestyle{plain}
\include{chapters/introduction}
}

%%% ČÁST I — Teoreticko-analytická
\part{Teoreticko-analytická část}
\chapter{Integrační architektury a jejich evoluce}
\label{ch:integracni-architektury}

\section{Výměna dat mezi systémy}

\section{Problém N×M integrace}

\section{Limity RAG a uzavřených systémů}

\section{Rešerše: existující MCP servery a standardizace ve vzdělávání}

\chapter{Model Context Protocol}
\label{ch:mcp}

\section{Původ a účel protokolu}

\section{Architektura klient–server a transport}

\section{Integrační primitiva}

\subsection{Nástroje (Tools)}

\subsection{Zdroje (Resources)}

\subsection{Prompty (Prompts)}

\section{Bezpečnostní model a rizika}

\chapter{Kontextové okno, tokeny a progressive disclosure}
\label{ch:kontext-tokeny}

\section{Kontextové okno a tokenizace}

\section{Náklady a latence předávání kontextu}

\section{Progressive disclosure vs. hromadné dumpování}

\chapter{Analýza domény a požadavků}
\label{ch:analyza}

\section{Doména studijní agendy VŠE a systém InSIS}

\section{Zdroj dat: Oracle isis-db2}

\section{Rozsah dat o vyučujících a GDPR}

\section{Funkční požadavky a případy užití}

\section{Metodika evaluace}


%%% ČÁST II — Praktická
\part{Praktická část}
\chapter{Návrh architektury}
\label{ch:navrh}

\section{Mapování entit InSIS na primitiva MCP}

\section{Přístup k datům bez ETL: přímé čtení z Oracle}

\section{Návrh nástrojů a progressive disclosure}

\section{Provozní architektura a nasazení}

\section{Slepé uličky a nečekané problémy}

\chapter{Implementace}
\label{ch:implementace}

\section{Jádro MCP serveru}

\section{Implementace nástrojů}

\section{Bezpečnost}

\section{Testování}

\section{Kontejnerizace, nasazení a CI/CD}

\section{Instrumentace a monitoring}

\chapter{Evaluace a diskuse}
\label{ch:evaluace}

\section{Datová propustnost a spotřeba tokenů}

\section{Latence a odezva}

\section{Matice splnění funkčních požadavků}

\section{Diskuse: vhodnost MCP pro univerzitní IS}

\section{Doporučení pro další výzkum}


%%% ZÁVĚR (nečíslovaná kapitola)
{%
\pagestyle{plain}
\include{chapters/conclusion}
}

%%% Seznam použité literatury
%%% Bibliography
%% Samostatná bibliografická databáze (references.bib)
\printbibliography[title={\bibnamex},heading={bibintoc}]


\end{document}
